\chapter{Insomnia}

\section*{Definition}
Insomnia is difficulty initiating sleep, maintaining sleep, or waking too early despite adequate opportunity for sleep, with daytime impairment or distress.

\section*{When to See a Doctor}
See your doctor if:
\begin{itemize}
  \item Persistent insomnia that affects daytime function, causes safety concerns, or occurs with snoring, breathing pauses, mania, depression, or suicidal thoughts
\end{itemize}

\section*{How It Develops}
Insomnia may be short-term or chronic. It can be maintained by conditioned arousal, irregular sleep schedules, stress, substances, medicines, medical illness, or another sleep disorder.

\section*{Causes}
\begin{itemize}
  \item Stress/anxiety
  \item Depression
  \item Circadian rhythm disruption
  \item Medical conditions (pain GERD sleep apnea)
  \item Medication/stimulant effects
\end{itemize}

\section*{Related Symptoms}
\begin{itemize}
  \item Daytime fatigue
  \item Irritability
  \item Concentration difficulty
  \item Mood changes
  \item Decreased performance
\end{itemize}


\section*{Medical Evaluation}
\begin{itemize}
  \item History assesses sleep schedule, time in bed, naps, substances, medicines, mood, pain, breathing symptoms, and daytime sleepiness.
  \item A sleep diary or validated insomnia questionnaire may help; examination and tests are directed at suspected medical or sleep disorders.
  \item Sleep studies are considered when sleep apnea, periodic limb movements, narcolepsy, or another sleep disorder is suspected.
\end{itemize}

\section*{Treatment Options}
\begin{itemize}
  \item Cognitive behavioral therapy for insomnia (CBT-I) is the recommended first-line treatment for chronic insomnia.
  \item Medicines may be considered short-term or in selected cases after reviewing benefits, risks, interactions, and causes; sedatives should not be self-started.
  \item Treat coexisting medical, psychiatric, circadian, or sleep disorders and review medicines that worsen sleep.
\end{itemize}

\section*{Self-Care and Home Management}
\begin{itemize}
  \item Establish a consistent daily routine with regular sleep, meals, and physical activity to support mental health stability.
  \item Practice stress management techniques including mindfulness, deep breathing exercises, and progressive muscle relaxation.
  \item Maintain social connections and avoid isolation by reaching out to trusted friends, family, or support groups.
  \item Avoid alcohol, caffeine, and recreational drugs which can destabilize mood and exacerbate psychiatric symptoms.
  \item If experiencing suicidal thoughts or crisis, contact emergency services or a crisis helpline immediately.
\end{itemize}

\section*{References}

\begin{thebibliography}{99}
\bibitem{insomnia:icsd3-insomnia} American Academy of Sleep Medicine. \emph{International Classification of Sleep Disorders}. 3rd ed, Text Revision. AASM; 2023. Chronic Insomnia Disorder.
\bibitem{insomnia:harrison-insomnia} Longo DL, et al. \emph{Harrison's Principles of Internal Medicine}. 21st ed. McGraw-Hill; 2022. Chapter 27: Insomnia.
\bibitem{insomnia:uptodate-insomnia} Manber R, et al. Evaluation and diagnosis of insomnia in adults. In: UpToDate, Post TW (Ed), UpToDate, Waltham, MA; 2025.
\bibitem{insomnia:morin} Morin CM, et al. Cognitive behavioral therapy for insomnia: a clinical practice guideline. \emph{JAMA}. 2023;329(9):739-749.
\bibitem{insomnia:buysse} Buysse DJ, et al. Insomnia: a clinical review. \emph{Lancet}. 2023;401(10378):813-826.
\bibitem{insomnia:kryger-insomnia} Kryger MH, et al. \emph{Principles and Practice of Sleep Medicine}. 7th ed. Elsevier; 2022. Chapter 70: Insomnia.
\end{thebibliography}
